\section{Dados e Métodos}
\label{sec:methods}

\subsection{Conjunto de Dados e Preprocessamento}
\label{sec:dataset}

Utiliza-se um conjunto de dados de recortes multiespectrais com dimensão
$64\times64$ pixels e quatro bandas (B2, B3, B4 e B8) provenientes de quatro
sensores: Sentinel-2, Landsat-8, Landsat-9 e MODIS. As estatísticas do conjunto
de dados, incluindo número de recortes por sensor e partição em treino,
validação e teste, são apresentadas na \Cref{tab:dataset-stats}.

As lacunas são introduzidas por máscaras sintéticas binárias (1 indica pixel
ausente; 0 indica pixel válido), com fração de lacuna aproximada de 10\%.
Adicionalmente, considera-se degradação por ruído em quatro níveis
(\texttt{inf}, 40, 30 e 20 dB), o que permite avaliar robustez sob cenários
controlados. Em todas as análises, as métricas de qualidade são calculadas
estritamente sobre os pixels mascarados, o que isola a contribuição do método
de preenchimento.

\begin{table*}[t]
 \centering
 \caption{Estatística do conjunto de dados por sensor de satélite.
 Todos os patches são de $64 \times 64$ pixels com 4 bandas espectrais (vermelha, azul, verde e IR).
 Divisão dos dados: 80\% treino, 10\% validação e 10\% teste.}
 \label{tab:dataset-stats}
 \begin{tabular}{lrrrrrr}
 \toprule
 Sensor & Banda & Resolução & \#Patches & \#Patches (teste) & \#Patches (Validação)  & Fração de Lacuna\\
 \midrule
Sentinel-2 & 4 & 10\,m & 73\,984 & 7\,399 & 7\,398 & $\approx$10\% \\
Landsat-8 & 4 & 30\,m & 1\,936 & 195 & 193 & $\approx$10\% \\
Landsat-9 & 4 & 30\,m & 1\,936 & 195 & 193 & $\approx$10\% \\
MODIS & 4 & 500\,m & 60 & 6 & 6 & $\approx$10\% \\
\textbf{Total} & -- & -- & \textbf{77\,916} & \textbf{7\,795} & \textbf{7\,790} & $\approx$10\% \\
 \bottomrule
 \end{tabular}
\end{table*}


\subsection{Entropia Local}
Seja uma variável aleatória discreta com distribuição $p(x)$ estimada a partir
do histograma local de intensidades. A entropia de Shannon é definida como
$H = -\sum_x p(x)\log p(x)$. Neste trabalho, a entropia local é calculada
pixel a pixel, em uma janela quadrada deslizante de lado $s\in\{7,15,31\}$,
utilizando uma transformação para tons de cinza obtida pela média das bandas.
Antes do cálculo, os valores são reescalados para o intervalo $[0,255]$ e
quantizados em \texttt{uint8}, o que viabiliza um filtro de ordem baseado em
janela para estimativa eficiente.

Como a interpretação de $H$ depende da escala espacial, cada janela é tratada
como um preditor distinto nas análises, evitando colapsar múltiplas escalas em
um único resumo.

\subsection{Métodos Clássicos}
Avaliam-se 15 métodos clássicos agrupados em seis famílias: interpolação
espacial, métodos baseados em kernel, krigagem geoestatística, inpainting no
domínio de transformadas e regularização, sensoriamento compressivo e métodos
patch-based. A \Cref{tab:methods} sumariza os métodos e seus parâmetros.

\begin{table}[htbp]
\centering
\caption{Visão geral dos 15 métodos clássicos de preenchimento de lacunas avaliados.}
\label{tab:methods}
\footnotesize
\resizebox{\columnwidth}{!}{%
\begin{tabular}{llp{3cm}}
\toprule
Categoria & Método & Parâmetros \\
\midrule
Espacial & Nearest Neighbor & -- \\
 & Bilinear & -- \\
 & Bicubic & -- \\
 & Lanczos (PG) & a=3 \\
\midrule
Kernel & IDW & p=2 \\
 & RBF & kernel TPS \\
 & Thin Plate Spline & -- \\
\midrule
Geoestatístico & Ordinary Kriging & variogram\_model=spherical, \newline nlags=6, \newline max\_points=2500, \newline random\_seed=13, \newline kernel\_size=None \\
\midrule
Transformada & DCT & max\_iterations=50, $\lambda$=0.05 \\
 & Wavelet & wavelet=db4, \newline level=3, \newline max\_iterations=50, \newline $\lambda$=0.05 \\
 & Total Variation & iter=100 \\
\midrule
Compressivo & L1-DCT (CS) & -- \\
 & L1-Wavelet (CS) & -- \\
\midrule
Baseado em Recortes & Non-Local Means & h=0.8, p=5, s=6 \\
 & Exemplar-Based & -- \\
\bottomrule
\end{tabular}
}
\end{table}


\subsection{Baselines de Aprendizado Profundo}
Cinco baselines de aprendizado profundo (\gls{AE}, \gls{VAE}, \gls{GAN}, U-Net e
\gls{ViT}) são incluídos para comparação no sensor Sentinel-2. Esses modelos
fornecem um ponto de referência contemporâneo, sem pretensão de exaustividade
arquitetural. A \Cref{tab:dl-architectures} resume as arquiteturas e funções de
perda consideradas.

\begin{table}[t]
 \centering
 \footnotesize
 \caption{Resumo das arquiteturas de \gls{DL}. Entrada:
 $[\mathbf{x};\, \mathbf{z}] \in \mathbb{R}^{5 \times 64 \times 64}$;
 Saída: $\hat{\mathbf{y}} \in [0,1]^{4 \times 64 \times 64}$.}
 \label{tab:dl-architectures}
 \resizebox{\columnwidth}{!}{%
 \begin{tabular}{llllc}
 \toprule
 Modelo & \emph{Bottleneck} / \emph{Core} & Função de Perda & \emph{Skip} & Otimizador \\
 \midrule
 Autoencoder (AE) & $512{\times}1{\times}1$ vector & MSE$_{\mathcal{G}}$ & -- & Adam \\
 Variational Autoencoder (VAE) & $\boldsymbol{\mu},\log\boldsymbol{\sigma}^2 \!\in\! \mathbb{R}^{256}$ & MSE$_{\mathcal{G}}$+$\beta$KL & -- & Adam \\
 Generative Adversarial Network (GAN) & Dilated conv (2,\,4) & ${\ell_1}^{\mathcal{G}}$+BCE$_{\mathrm{adv}}$ & \checkmark & Adam $(\beta_1{=}0.5,\beta_2{=}0.999)$ \\
 U-Net & $1024{\times}4{\times}4$ ResBlock & MSE$_{\mathcal{G}}$ & \checkmark & AdamW \\
 Vision Transformer (ViT) & 4$\times$ Transformer ($d$=256) & MSE$_{\mathcal{G}}$ & -- & AdamW \\
 \bottomrule
 \end{tabular}
 }
\end{table}

